% Section 1: Introduction (about 1 page). Plan: PAPER_PLAN 4 / §1.
%   p1 hook (LBSN, anticipate the next move, mobility-aware motivation).
%   p2 the two tasks (next-category, next-region; we do NOT predict the exact next place).
%   p3 the tension (sharing is not free; where it helps, where it costs).
%   p4 what we do (check-in-level representation + one model, two tasks).
%   p5 contributions (three bullets, stated positively).
%   p6 roadmap.
% Reminder: next-category / next-region; American English; no em-dash.
\section{Introduction}
\label{sec:intro}

Location-based social networks, such as Gowalla~\cite{cho2011gowalla}, let people share the places that they visit as check-ins, which together record how people move through a city. Individual mobility is highly predictable in principle~\cite{song2010limits}, so a service that anticipates the next move can prepare ahead of time, caching content where a user is heading~\cite{bastug2014edge} or provisioning capacity at the destination before demand arrives. Handover predictions already let cellular services adapt in advance at the network level~\cite{vielhaus2022handover}; at the city scale, check-in mobility analysis is likewise a design input for mobility-aware services~\cite{moura2025mobilityaware}.

Predicting the next place exactly, a single point of interest (POI), is hard and often more than a service needs. Two coarser questions are usually enough: the next category, the type of place such as food or shopping, and the next region, the part of the city that a user is heading to (a neighborhood-scale unit such as the U.S. census tract). The first captures intent, the second captures geography; we study whether one model should learn both at once.

Sharing one representation across tasks has a cost: in multi-task learning (MTL), one model does several jobs at once by sharing most of its parts, so the shared parameters can converge to a compromise optimal for neither task, helping one while hurting the other~\cite{caruana1997multitask}. Our earlier
work reported no consistent multi-task advantage for the paired category tasks and attributed it, in part, to this effect~\cite{silva2025mtlnet}; the useful question is where sharing helps, what it costs, and how to share so the gains hold and the cost stays small.

We propose two enhancements. First, we build a check-in-level representation: instead of one fixed vector per place, each check-in gets its own vector that holds its context (the time, nearby places, and recent visits).
This adds a fourth level, the check-in, beneath the place, region, and city levels of hierarchical graph infomax~\cite{huang2023hgi,velickovic2019dgi} (Fig.~\ref{fig:dataflow}). Second, we train one model that
predicts the next category and the next region in a single forward pass, with a shared trunk (a cross-attention stack where the two tasks exchange
semantic context) and a private spatial path for the region task (Fig.~\ref{fig:model}). We evaluate on two settings chosen to differ: five U.S. states of different sizes (Gowalla) and one non-U.S. city (Istanbul).\footnote{\label{fn:code}Code (model, representation, baselines, statistical tests): \url{https://github.com/VitorHugoOli/PoiMtlNet/tree/mobiwac}. Both data sources are public: the category-annotated Gowalla dump (\url{https://figshare.com/articles/dataset/gowalla_data/22126586}) and Massive-STEPS~\cite{wongso2025massivesteps}.}

Our contributions are as follows.
\begin{itemize}
  \item A check-in-level representation that extends hierarchical graph infomax from the place to the individual visit, improving next-category prediction over a standard place embedding at every one of the six datasets, by $+0.2$ to $+6.3$ points of macro-F1 (a balanced score that counts every
    category equally), with every fold favoring it at every dataset (Table~\ref{tab:substrate}). The advantage tracks the geometry of the per-visit
    vectors, which separate categories at the visit level and not at the place level. A single fixed vector
    per place cannot separate places that serve several purposes, and
    per-visit context recovers that distinction.
  \item A single model for joint next-category and next-region prediction, sharing semantic context across tasks while keeping a private spatial path for the region task; to our knowledge, the first work to treat fine-grained region as an end target of equal standing (Section~\ref{sec:related-tasks}). In one forward pass, and reading both answers from a single saved model, it outperforms the dedicated
    region model at the two datasets with the most regions (by about a point) and stays
    statistically non-inferior within a two-point margin (TOST, a test of
    equivalence) at the other four. On category it outperforms the dedicated model at one
    dataset and stays within a fifth of a point of it at the rest, even though each
    dedicated model is tuned per dataset (Table~\ref{tab:results}).
  \item An empirical account, across five U.S. states and one non-U.S. city, of where joint training pays.
    The two datasets where the joint model outperforms the dedicated region model are the two with the
    largest region vocabularies (Texas and California), which is where the region task is hardest;
    at the four smaller vocabularies the two approaches are equivalent within the margin
    (Fig.~\ref{fig:deltas}). The gain does not order monotonically with region count inside that
    pair, so we report the grouping rather than a law. All joint and dedicated results average
    four random initializations (Section~\ref{sec:results-part2}).
\end{itemize}

The rest of the paper presents related work (Section~\ref{sec:related}), the problem (Section~\ref{sec:problem}), the method and setup (Sections~\ref{sec:method} and~\ref{sec:setup}), results (Section~\ref{sec:results}), and discussion and conclusions (Sections~\ref{sec:discussion} and~\ref{sec:conclusion}).
